\section{Seven lifts from a compatible line class}\label{app:p48}

The antipodal block lift amplifies improvements to a single compatible
class. Two choices determine the gain: the size of the class and the
overlap among its selected images. This section makes both choices
explicit for $P_{48p}$ and obtains seven configurations from one
geometric argument.

\subsection{The tail configurations}
For dimensions $k=1,\ldots,7$, use respectively the normalized root
configurations
$A_1,A_2,D_3,D_4,D_5,E_6,E_7$. They have
$2,6,12,24,40,72,126$ points and are antipodal. Explicitly,
$A_1=\{-1,1\}$, and $A_2$ is the six points
$(e_i-e_j)/\sqrt2$ in the plane $x_1+x_2+x_3=0$. The $D_k$ roots
are $(\eps e_i+\eps'e_j)/\sqrt2$, with $i<j$. Their kissing
conditions follow directly from their coordinate overlaps.

For completeness the last two codes can be specified together in
an integer model. Let
\begin{equation}\label{eq:E8}
 W=\{\pm2e_i\pm2e_j:i<j\}
  \cup\{(\eps_0,\ldots,\eps_7):\eps_i\in\{\pm1\},\ \prod_i\eps_i=1\}.
\end{equation}
Every row has squared norm eight and distinct rows have inner product
at most four: this follows from support overlap for the first family,
at most two signed contributions for cross-family pairs, and an even
positive sign distance within the second family. The subsets
\[
 W_7=\{w\in W:w_6+w_7=0\},\qquad
 W_6=\{w\in W:w_5=w_6=-w_7\}
\]
have respectively $62+64=126$ and $40+32=72$ rows. The maps
\[
 w\longmapsto(w_0,\ldots,w_5,\sqrt2w_6),\qquad
 w\longmapsto(w_0,\ldots,w_4,\sqrt3w_5)
\]
are isometries on the indicated subspaces. Dividing by $\sqrt8$
therefore gives the required $7$- and $6$-dimensional antipodal codes.


\subsection{The 7077-line class and its automorphism images}
\label{sec:p48-update}
We use the catalogue lattice $P_{48p}$ in its published
integral Gram basis $G$~\cite{p48p-catalogue}. Its minimum squared norm
is six. Since it is even unimodular of rank 48, the theta-series
calculation in~\eqref{eq:theta48} gives 52416000 minimal vectors.
The witnesses use this catalogue basis directly. The quadratic-residue
neighbour used in dimension 45 is constructed in
Section~\ref{sec:qr-neighbour}. The theta-series calculation in
Appendix~\ref{app:pless-neighbour} applies to every extremal even
unimodular lattice of rank 48.

Represent a minimal line by an integer row $c$ satisfying
$cGc^t=6$, identifying $c$ and $-c$. The new mother class
$\mathcal C$ has 7077 lines and satisfies
\begin{equation}\label{eq:p48-class}
 |cGc'^t|\le2\qquad(c\ne\pm c',\ c,c'\in\mathcal C).
\end{equation}
Every norm and pair product in this assertion is an integer.
The certificate supplies the complete class and matrices $A_i$ with
$A_iGA_i^t=G$. Thus $\mathcal C A_i$ is another class with the
same property.

For $t$ selected images, assign each line to the first image in which
it occurs. Write $\mathcal D_i$ for the assigned subset of
$\mathcal C A_i$, and put
\[
 S=\left|\bigcup_{i=1}^t\mathcal C A_i\right|
   =\sum_{i=1}^t|\mathcal D_i|.
\]
The $\mathcal D_i$ are disjoint as line sets. Taking both signs of
each assigned line gives disjoint vector blocks $B_i$ of size
$2|\mathcal D_i|$, with normalized inner products at most $1/3$
by~\eqref{eq:p48-class}. Assign these blocks injectively to the $t$
distinct antipodal lines of the corresponding root tail code.
Theorem~\ref{thm:block-lift} gives
\begin{equation}\label{eq:p48-lines-count}
 K(48+k)\ge52416000+2S+\tau_k,
 \qquad \tau_k=2t.
\end{equation}

The injective assignment is part of the construction. For two
different assigned classes, the minimal-norm bound gives a normalized
head product at most $1/2$, and distinct tail lines give a tail
product at most $1/2$. Their cap product is therefore at most
$(3/4)(1/2)+(1/4)(1/2)=1/2$. Reusing a tail line for arbitrary
different classes would instead give $5/8$ and is not permitted.

\begin{theorem}\label{thm:p48-update}
The class and image matrices in the supplementary data give the
seven bounds in Table~\ref{tab:p48-update}.
\end{theorem}
\begin{proof}
The finite checker verifies the catalogue Gram matrix, its exact
isometries, all 7077 norms, distinctness modulo sign, and
all distinct-line inequalities~\eqref{eq:p48-class}. It then
reconstructs each selected image and counts its union with the
preceding images. This gives the $S$ values in the table.
The root tail codes have the ranks, sizes, and inner products
established above. Hence~\eqref{eq:p48-lines-count} applies.
\end{proof}

\begin{table}[ht]
\centering
\caption{The line-class liftings. The union, rather than
the sum of image sizes, determines the number of available lines.}
\label{tab:p48-update}
\begin{tabular}{rrrrr}
\toprule
$d$&$t$&$S$&$\tau_k$&$52416000+2S+\tau_k$\\\midrule
49&1&7077&2&52430156\\
50&3&21231&6&52458468\\
51&6&42459&12&52500930\\
52&12&84874&24&52585772\\
53&20&141328&40&52698696\\
54&36&253851&72&52923774\\
55&63&442507&126&53301140\\\bottomrule
\end{tabular}
\end{table}

\subsection{Class growth and image selection}
The public comparison uses a 7069-line class
$\mathcal C_0\subset\mathcal C$~\cite{kissingnumbers}. With the selected
matrices held fixed, let
\[
 S_0=\left|\bigcup_{i=1}^t\mathcal C_0 A_i\right|,
\]
and let $S_{\rm ref}$ be the union size in the comparison construction.
The increase in the final point count splits exactly as
\begin{equation}\label{eq:p48-gain}
 2(S-S_{\rm ref})=2(S-S_0)+2(S_0-S_{\rm ref}).
\end{equation}
The first term measures the effect of enlarging the class with these
images fixed; the second measures the effect of changing the images
with the original class fixed. Both are obtained by exact union counts.

\begin{table}[ht]
\centering
\caption{The two contributions to the seven point-count improvements.}
\label{tab:p48-gain}
\begin{tabular}{rrrrr}
\toprule
$d$&$S_0$&$2(S-S_0)$&$2(S_0-S_{\rm ref})$&Total\\\midrule
49&7069&16&0&16\\
50&21207&48&2&50\\
51&42411&96&18&114\\
52&84778&192&64&256\\
53&141168&320&154&474\\
54&253565&572&296&868\\
55&442009&996&414&1410\\\bottomrule
\end{tabular}
\end{table}

The three selected images in dimension 50 are disjoint. Larger image
families overlap, so an eight-line enlargement need not supply eight
new lines per image. In dimension 55 it contributes 498 distinct lines,
rather than $8\cdot63=504$. The image choice contributes a further
207 lines. Thus the two improvements reinforce one another through
the union, rather than through a sum of independently optimized class sizes.
